\section{Spatial comparison at a vanishing field}\label{sec:weak-field}
\label{sec:spatial-comparison}

The sparse-pin estimate controls the missing cycles at prescribed
points.  We now turn it into a bound for a spatially oscillating
magnetization.  The two steps are quite different.  First, an exact
susceptibility identity expresses the transverse square through a
one-particle resolvent on the holes.  Second, we show that the resolvent
is nearly constant on large rectangular cells.  Small closed classes of
holes are allowed; a trace identity controls their total contribution.

\begin{proposition}[Weak-field spatial estimate]\label{prop:weak-field}
Fix $S$.  For all sufficiently large $\beta$ there are
$C_\beta,R_0(\beta)<\infty$ with the following property.  Let $L$ be
even, let $R_0(\beta)\le R\le L$, and partition $\Lambda_L$ into
rectangular cells whose side lengths belong to $[R/2,4R]$.  For every
deterministic real $f:\Lambda_L\to[-1,1]$ with
$\sum_{x\in B}f_x=0$ on each cell $B$, put
\[
 D_f=\sum_x f_x Z_x,\qquad q=R^{-5/2},\qquad
 w=1-q=e^{-b}.
\]
If $\langle\cdot\rangle_b$ is the Gibbs expectation with exponent
$H_0+b\sum_xY_x$, then
\begin{equation}\label{eq:weak-field}
 \frac qV\langle D_f^2\rangle_b\le C_\beta R^{-1/40}.
\end{equation}
The constants are uniform in the partition, $f$, and $L$.
\end{proposition}

Throughout the proof $\beta$ is fixed above a sufficiently large
threshold; subsequent lower bounds on $R$ may depend on $\beta$.
We use the soft-pin representation of Section~\ref{sec:sparse-pins},
with no extra deterministic real pins.  Expose all cycles missing the
soft pins.  Write $\mathbf a$ for their trajectories, $A_0$ for their
real slots at the cut, and $\mathcal H$ for the complementary cut
slots, called the holes.  The cut is after the seam and its soft-pin
insertions, which do not move any line.
The soft pins themselves may be placed on real cut slots, instead of
using the equivalent ghosts.  A hole cycle is then pinned and has up
color; every exposed cycle receives an independent fair color.

\subsection{The susceptibility as a hole resolvent}

Condition on $\mathbf a$.  The fresh real-hole motion consists of
rate-$1/2$ jumps across the currently available slot edges, the forced
hole transports at exposed jumps, and a uniform completion of the seam
permutation within each site.  Its one-turn kernel, without killing
or conditioning on its cycles, is denoted by $K$.  In this section we
use the \emph{row-forward convention}: $K_{ij}$ is the probability
of going from $i$ to $j$.  Thus $K$ is doubly stochastic.  Set
\begin{equation}\label{eq:hole-resolvent}
 D=I-wK,\qquad D^{-1}=\sum_{r\ge0}w^rK^r.
\end{equation}
This inverse exists and is entrywise nonnegative.  Empty hole sets use
empty matrices, determinant one, and sums zero.

\begin{lemma}[Soft-pin identities]\label{lem:soft-resolvent}
Given the exposure $\mathbf a$, for $i,j\in\mathcal H$, $i\ne j$,
\begin{align}
 \PP(i\text{ is pinned}\mid\mathbf a)
   &=q(D^{-1})_{ii},\label{eq:pin-marginal}\\
 \PP(i,j\text{ are pinned, next pin from }i\text{ is }j
            \mid\mathbf a)
   &=\frac{q^2}{w}(D^{-1})_{ij}.\label{eq:pin-successor}
\end{align}
Moreover,
\begin{equation}\label{eq:resolvent-trace}
 \E\Tr D^{-1}\le N.
\end{equation}
For $f_i=f_{x_i}$ on the slots,
\begin{equation}\label{eq:transverse-resolvent}
 \langle D_f^2\rangle_b
 =\frac14\sum_i f_i^2+
   \frac12\E\sum_{\substack{i,j\in\mathcal H\\i\ne j}}
                   f_if_j(D^{-1})_{ij}.
\end{equation}
All expectations here include the normalized picture and soft-pin law.
\end{lemma}

\begin{proof}
In the exposure disintegration, the proposed exposed paths have a
fixed cycle weight and a factor $w^{|A_0|}$, since none of their cut
slots is pinned.  For a proposed pin set $P'\subseteq\mathcal H$,
put $Q=\mathcal H\setminus P'$.  Lemma~\ref{lem:next-pin}, with
the row convention obtained by transposition, says that the remaining
cycles all meet $P'$ with probability $\det(I-K_{QQ})$.  Consequently
the unnormalized conditional weight of $P'$ is
\begin{equation}\label{eq:conditional-soft-weight}
 q^{|P'|}w^{|\mathcal H|-|P'|}\det(I-K_{QQ}).
\end{equation}
Multilinearity of the determinant in the rows gives
\[
 \sum_{P'\subseteq\mathcal H}
 q^{|P'|}w^{|\mathcal H|-|P'|}\det(I-K_{QQ})=\det(I-wK).
\]
Restricting this sum to sets containing $i$ gives $q$ times the
diagonal cofactor of $D$, proving \eqref{eq:pin-marginal}.

For \eqref{eq:pin-successor}, keep $i,j$ pinned.  At fixed $Q$ the
all-hit probability times the conditional next-pin entry is
\[
 \det\begin{pmatrix}
 K_{ij}&K_{iQ}\\ -K_{Qj}&I-K_{QQ}
 \end{pmatrix}.
\]
This is a polynomial identity, including when $I-K_{QQ}$ is singular,
by the unnormalized identity in Lemma~\ref{lem:next-pin}.
Put $G=\mathcal H\setminus\{i,j\}$.  Summing over the Bernoulli
choices on $G$ replaces this expression by
\[
 \det D_{GG}\bigl(K_{ij}+wK_{iG}D_{GG}^{-1}K_{Gj}\bigr).
\]
Indeed an unselected row in the bordered determinant is an identity
row with zero border entry, and a selected row has probability $w$.
The principal submatrix $D_{GG}$ is invertible because $w<1$.
The two-by-two Schur complement of $D_{GG}$ in $D$ now gives the
last display as $\det D\,(D^{-1})_{ij}/w$.  Multiply by $q^2$
and divide by the normalizing determinant to obtain the assertion.

Condition instead on the picture and its colors.  Each up cut slot
is independently pinned with probability $q$, and no down slot can
be pinned: summing its two allowed Bernoulli weights gives $w+q=1$
for up, and only $w$ for down.  The expected total number of pins
is therefore at most $qN$.  Pins never lie in $A_0$, so summing
\eqref{eq:pin-marginal} and taking expectations proves
\eqref{eq:resolvent-trace}.

For the transverse identity rotate coordinates so that the field
component is $s^3$ and the tested component is $s^1$.  The normalized
field matrix is $\diag(1,w)$.  Split it symmetrically around each
insertion; this is permitted because the total field commutes with
$H_0$.  At a transverse insertion its slot matrix is
\[
 \diag(1,\sqrt w)\,s^1\,\diag(1,\sqrt w)
   =\frac{\sqrt w}{2}\begin{pmatrix}0&1\\1&0\end{pmatrix}.
\]
For distinct $i,j$, expand pins on all other slots.  Two flips give
zero unless they lie on the same cycle.  On that cycle one of the two
arcs between the flips must be down and must contain no other pin.
Adding pins at $i,j$ therefore identifies each allowed coloring with
one of the events ``next pin from $i$ is $j$'' or ``next pin from
$j$ is $i$.''  If there are no other pins on this cycle both events
contribute, exactly accounting for its two allowed colorings.
The ratio of insertion weights to the weights with the added pins is
$w/(4q^2)$.  Formula~\eqref{eq:pin-successor} thus gives
$\frac14\E[(D^{-1})_{ij}+(D^{-1})_{ji}]$ for the two-slot
expectation.  Sum against $f_if_j$ over ordered distinct pairs.
Here inverse entries are interpreted as zero if either cut slot is
outside the random set $\mathcal H$.
The same-slot identity $(s^1)^2=I/4$ gives the first term in
\eqref{eq:transverse-resolvent}.  The seam projection ensures that
these slot sums are precisely the physical site-spin observable.
\end{proof}

Define the bounded real vector
\begin{equation}\label{eq:resolvent-u}
 u=qD^{-1}f_{\mathcal H},\qquad
 e_i=\sum_{j\in\mathcal H}K_{ij}(u_i-u_j)^2.
\end{equation}
Since $qD^{-1}$ is stochastic, $|u_i|\le1$.  Double stochasticity
and $Du=qf_{\mathcal H}$ imply
\begin{equation}\label{eq:traversal-energy}
 \frac12\sum_i e_i
 =\ip{u}{(I-K)u}
 =\frac q w\bigl(\ip{u}{f_{\mathcal H}}-\norm{u}_2^2\bigr)
 \le CqN.
\end{equation}
No reversibility of $K$ is used here.

We will approximate $u$ by a constant on the holes in each cell.
The missing slots mean that $f_{\mathcal H}$ no longer has exact
cellwise cancellation.  The following estimate supplies its replacement.

\begin{lemma}[Cancellation after exposure]\label{lem:hole-cancellation}
For every cell $B$,
\begin{equation}\label{eq:hole-cancellation}
 \E\left|\sum_{\substack{i\in\mathcal H\\x_i\in B}}f_i\right|
 \le C\sqrt{|B|/q}.
\end{equation}
\end{lemma}

\begin{proof}
Return to the original axes, with field direction $Y$, and put
$T_B=\sum_{x\in B}f_xY_x$.  The ordinary Laplace transform
$F_B(z)=\langle e^{zT_B}\rangle_b$ is zero-free on $|z|<cb$:
all slot fields in direction $Y$ have positive real part there,
so Lemma~\ref{lem:trace-zero-free} applies.  Take its analytic
logarithm equal to zero at $0$.  Spectral calculus for $T_B$ gives
\[
 \Rea\log F_B(z)\le |z|S|B|\le Cb|B|\qquad(|z|<cb).
\]
The Borel--Carath\'eodory inequality on this disk, followed by
Cauchy's second-derivative estimate on a smaller concentric disk,
therefore yields
\[
 \Var_b(T_B)=(\log F_B)''(0)\le C|B|/b\le C|B|/q.
\]
Translation invariance and $\sum_{x\in B}f_x=0$ give
$\langle T_B\rangle_b=0$.  In the color representation its
conditional expectation given $\mathbf a$ is
\[
 \E(T_B\mid\mathbf a)
 =-\frac12\sum_{\substack{i\in A_0\\x_i\in B}}f_i
 =\frac12\sum_{\substack{i\in\mathcal H\\x_i\in B}}f_i.
\]
Here exposed cycles are fair and every hole cycle is up; the total
sum of the slot coefficients on $B$ is zero.  Conditional Jensen
and Cauchy--Schwarz prove \eqref{eq:hole-cancellation}.
\end{proof}

\subsection{A graph recording communication between holes}

The small traversal energy in \eqref{eq:traversal-energy} must now
control spatial variation.  A hole can be carried by the exposed
trajectories even when the fresh walk makes no jumps.  We retain
that motion explicitly in a graph on the cut holes.

Ignore fresh walk jumps and follow only forced transports from the
cut to the seam.  Choose one allowed residual seam matching within
each site.  The resulting lines, called \emph{skeleton lines}, give
a permutation $\phi$ of $\mathcal H$.  Define an undirected graph
$\Gamma$ on $\mathcal H$ with the following edges:
\begin{enumerate}[label=(\roman*),nosep]
 \item an endpoint edge from $i$ to every possible final cut hole
 under a residual seam matching, including $\phi(i)$;
 \item a contact edge between $i,j$ whenever their skeleton lines
 occupy adjacent sites, with an available slot edge between them,
 throughout some interval of positive length.
\end{enumerate}
Self-edges may be omitted.  All choices of representatives and
matching can be made by fixed orderings of the finite sets.

\begin{lemma}[Communication costs and closed classes]
\label{lem:skeleton-graph}
An endpoint edge from $i$ to $j$ satisfies
$(u_i-u_j)^2\le C_\beta e_i$.
For a contact edge witnessed by a fixed-position interval of length
$\tau\le1$,
\begin{equation}\label{eq:contact-cost}
 (u_i-u_j)^2\le C_\beta\tau^{-1}(e_i+e_j).
\end{equation}
The kernel $K$ is closed on every connected component of $\Gamma$,
and
\begin{equation}\label{eq:component-count}
 \E\,\#\{\text{components of }\Gamma\}\le qN.
\end{equation}
\end{lemma}

\begin{proof}
The walk jump rate is uniformly bounded by a constant depending on
$\ell$.  The probability of making no fresh jump during a period
is at least $e^{-C\beta}$; the probability of any prescribed allowed
seam image is at least $1/\ell$.  Thus an endpoint entry of $K$ is
at least $c_\beta$.  For a contact interval, require exactly one
jump from line $i$ to line $j$ during that interval and no other
jumps, and choose endpoint $\phi(j)$.  This gives
$K_{i,\phi(j)}\ge c_\beta\tau$.  Compare $u_i$ to $u_j$
through $u_{\phi(j)}$, using also the endpoint bound for $j$;
this proves \eqref{eq:contact-cost}.

A fresh walk changes skeleton line only through contact, and every
possible seam endpoint has been included in $\Gamma$.  Hence each
component is closed under $K$.  Its stochastic block has eigenvalue
$1$.  Every eigenvalue $\lambda$ of that block satisfies
$|\lambda|\le1$, and
\[
 \Rea\frac1{1-w\lambda}
 =\frac{1-w\Rea\lambda}{|1-w\lambda|^2}>0.
\]
Its contribution to $\Tr D^{-1}$ is therefore at least $1/q$,
from the eigenvalue $1$.  This trace argument applies equally to
non-diagonalizable blocks.  Sum over the components and use
\eqref{eq:resolvent-trace}.
\end{proof}

Every edge of $\Gamma$ also has a spatial realization.  An endpoint
edge follows one skeleton line and then the seam, which does not
change its site.  A contact edge follows its first skeleton line to
a contact time, crosses one site edge, and follows the second line
back to its cut.  Such a realization consists of at most two pieces
of actual marks in chronological or reverse chronological order,
and one extra site edge.  This observation will prevent an edge
leaving a small obstacle from making a long excursion or winding
around the torus.

\subsection{Coarse obstacles and local communication}

Choose $d=\lceil C_0\beta\rceil$, with the fixed constant $C_0$
large enough below.  Partition each coordinate circle into integer
intervals of lengths in $[d,2d]$.  Their product boxes are indexed
by a three-dimensional coordinate torus.  Fine distances refer to
site coordinates and coarse distances to these box indices; both
use the supremum norm and the periodic convention.  A site is
\emph{blocked} if it has no cut hole.

We use the following deterministic input from
\cite[Lemma~5.5]{Ord}, also recalled in \cite[Section~3]{BL}.
For each fixed $a$ there is $\rho(3,a)\ge a+1$ such that two
unblocked sites at distance at most $a$ can be joined by a
nearest-neighbor unblocked path in the ball of radius $\rho(3,a)$
about the first, of length at most $(2\rho(3,a)+1)^3$, provided
that ball contains fewer than six blocked sites.  We use $a=44$
and fix $M\ge\max\{48,\rho(3,44)+2\}$.  The result concerns
arbitrary blocked sets; we apply it in local lifts to the lattice.
Increase $C_0$, if necessary, so that $d$ dominates these fixed
geometric constants.

Let $0<\delta<\min\{1,\beta/2\}$; its value in terms of $R$
will be chosen only after the comparison is proved.  Put
$j_* =\lfloor d/20\rfloor$.  Declare a coarse box \emph{bad}
when at least one of the following occurs:
\begin{enumerate}[label=(\roman*)]
 \item some $M$-ball centered within fine distance $10d$ of the
 box contains at least six blocked sites;
 \item starting within that $10d$-neighborhood, a site path of
 $j_*$ inter-site steps is supplied by distinct actual marks in
 chronological or reverse chronological order during one period;
 \item among marks on edges incident to the $15d$-neighborhood,
 there is a mark within $\delta$ of $0$ or $\beta$, or two distinct
 marks at times less than $\delta$ apart.
\end{enumerate}
The last event will be called a \emph{short-time flag}.  Slot
compatibility is deliberately not required in (ii); this makes the
condition constrain every skeleton trajectory.  Join bad boxes
at coarse distance at most eight.  All probabilities below use the
original joint picture law.  The bad events may depend on the full
picture, even though $K$ and $u$ depend only on its exposure.

\begin{lemma}[Bad-box tails]\label{lem:bad-box-tail}
There are $\beta_0,c,C,\lambda_1>0$ and, for each
$\beta\ge\beta_0$, a number $\delta_0(\beta)>0$ such that for
every $0<\delta\le\delta_0(\beta)$, uniformly in the admissible
volumes and pin parameters of this section,
\begin{equation}\label{eq:bad-box-tail}
 \PP\{z\text{ is bad and its bad component reaches distance }h
          \text{ from }z\}
 \le C e^{-ch^{\lambda_1}}.
\end{equation}
A specified box has a short-time flag with probability at most
$C_\beta\delta$.
\end{lemma}

\begin{proof}
Factorial domination from Lemma~\ref{lem:pin-inputs} bounds the
one-mark count near the ends of the period and the ordered two-mark
count with time gap less than $\delta$ by $C_\beta\delta$.
For a prescribed starting site, the expected number of ordered
paths using $j_*$ marks is at most $(C\beta)^{j_*}/j_*!$: the
bounded choices of a next slot edge and the factorial factor are
absorbed in $C$, and the ordered time simplex has volume
$\beta^{j_*}/j_*!$.  Summing starting sites gives the bound
\begin{equation}\label{eq:coarse-mark-cost}
 \varepsilon_{\rm mark}
 =Cd^3(C\beta)^{j_*}/j_*!+C_\beta\delta
\end{equation}
for the union of (ii) and (iii).  Such paths stay within fine
distance $11d$ of the box, so all their data are in its fixed
$15d$-neighborhood, with a harmless one-edge enlargement.
Take $C_0$ large; the first term is then at most $Ce^{-c\beta}$.

Here is the dependence estimate needed to pass from single boxes
to components.  Choose a large integer $a$ with $a^{1/100}>4$,
and a fixed coarse radius $h_0$ large enough that data neighborhoods
at separated leaves below are disjoint.  Set $h_k=h_0a^k$.
A bad path from a prescribed root to distance $h_k$ contains two
crossings at scale $h_{k-1}$: one starts at the root, and one at
its first visit to distance $h_k/2$.  Truncate each on its first
reach of the smaller radius.  Their ranges are separated by at
least $h_k/4$, since
\[
 h_k/2-2(h_{k-1}+8)\ge h_k/4.
\]
Recursion gives $2^k$ distinct bad leaves with the sparse growth
bound of Proposition~\ref{prop:sparse-pins}, in coarse coordinates.
Indeed a ball of radius less than $h_r/8$ meets at most one branch
at every split of scale at least $h_r$, and hence at most
$2^{r-1}$ leaves; the choice of $a$ gives the growth exponent
$1/100$ also between consecutive scales.  The number of possible
witnesses is at most $C_1^{2^k}$, because at a split of scale $h_r$
there are at most $Ch_r^3$ choices and
$\sum_{r\ge1}2^{-r}\log(Ch_r^3)<\infty$.

For any specified subset of leaves of type (i), choose a witnessing
center and six blocked sites.  There are at most $Cd^3$ choices per
leaf, since $M$ is fixed.  Select one fixed slot at each chosen
site.  These fine-slot tests are sparse with a constant independent
of $d$: sites in a fine ball of radius $r$ come from leaves in a
coarse ball of radius $C(1+r/d)$, so their number is at most
$6C[1+C(1+r/d)]^{1/100}\le C_s(1+r)^{1/100}$.
Leaf separation makes the selected sites distinct.  A blocked site
has all its slots exposed, so Proposition~\ref{prop:sparse-pins}
bounds the joint probability for $m$ such leaves by
$(Cd^3p^6)^m$, where $p=\beta^{-9/10}$.

For any subset of mark-type leaves, the data neighborhoods are
disjoint.  Expand their indicators by the path and flag counts
used above and apply factorial domination to all the selected marks
together.  Distinct neighborhoods use distinct marks; their
integrals factor.  This gives the product bound
$\varepsilon_{\rm mark}^m$, without assuming independence of the
weighted picture.  In every witness at least half the leaves have
one of these two types.  Sum over that subset and over the
witnesses.  Since $d^3p^6=O(\beta^{-12/5})$, choosing $\beta$
large and then $\delta_0(\beta)$ small puts both leaf costs below
a fixed sufficiently small witness budget.  The probability is then
at most $Ce^{-c2^k}$, with $c,C$ uniform for
$\delta\le\delta_0(\beta)$.  This proves \eqref{eq:bad-box-tail} with
$\lambda_1=\log_a2$, after changing constants for intermediate
and bounded radii.  The argument uses only wrapped distances and
the bound $|B_h|\le C(1+h)^3$, so it is valid on the coarse torus.
\end{proof}

\begin{lemma}[Local paths near good boxes]\label{lem:good-box-paths}
Choose a representative cut hole $i(z)$ in every good box $z$.
Any two cut holes in one good box, and representative holes in two
nearest-neighbor good boxes, can be joined in $\Gamma$ using at
most $C_\beta$ contact edges, all with contact duration at least
$\delta$, whose cut starts stay within fine distance $C_\beta$
of the boxes.

Moreover, if a spatial realization of any edge of $\Gamma$ meets
a good box $z$, its whole realization lies within fine distance
$d/2$ of a point in that box, after increasing the fixed lower
threshold on $d$.  In particular both cut endpoint boxes have
coarse distance at most one from $z$.
\end{lemma}

\begin{proof}
A good box contains an unblocked site, since otherwise an $M$-ball
in it would contain six blocked sites.  To connect two desired
sites, first use a coordinate path of length $C d$ through the
one- or two-box neighborhood.  Replace every blocked vertex of
this path by a surviving nearest neighbor.  Such a neighbor exists:
six blocked neighbors would contradict (i).  Consecutive replaced
vertices have distance at most three.  The imported local-connectivity
result joins them through unblocked sites with bounded length and
displacement.  All its balls lie in the $10d$-neighborhood where
(i) is absent.  Lift the resulting path to cut slots, using the
complete adjacency of available slots across a site edge.  If
different slots at one site need joining, pass through a surviving
neighbor of that site.  The absence of (iii) means that these cut
adjacencies persist for at least $\delta$ after the cut.  This
proves the first assertion.

For the second assertion, a realization has at most two monotone
mark pieces and a contact edge.  Suppose one piece meets $z$.
The portions on either side of that visit each have fewer than
$j_*$ steps: a longer portion, read in the appropriate time
direction, would violate (ii).  If there is a second piece, its
contact endpoint is therefore within $2j_*+1$ of the visit.
This point lies in the tested $10d$-neighborhood, and (ii) bounds
the whole second piece by $j_*-1$ steps.  If the contact edge
itself meets $z$, apply the same argument from its endpoints.
Thus every point of the realization lies within $3j_*+2<d/2$
of the visit.  This reasoning also applies in a lift: every
monotone piece uses distinct torus marks during one period.
\end{proof}

The first assertion converts ordinary paths through good boxes into
energy bounds.  The second is the essential control on a hole trapped
among bad boxes: any escape through a good boundary must take place
locally.  We next build those boundaries on a finite portion of the
covering lattice, including when the cell under consideration is the
whole torus.

\subsection{Filled hulls and the comparison function}

We recall precisely the deterministic hull construction that we
import from \cite[Lemma~6.1 and Section~6.1, Equations~(6.2)--(6.3)]{Ord}.
For any finite blocked set $\mathcal Z\subset\Z^3$, join its
points at supremum distance at most eight.  For each component $C$,
thicken it to $A_C=\{z:\dist(z,C)\le4\}$ and fill the finite
nearest-neighbor components of $\Z^3\setminus A_C$.  The resulting
hull $H_C$ is contained in the coordinate bounding box of $A_C$.
The hulls are disjoint or nested, so the maximal ones are disjoint.
Their exterior vertex boundaries consist of unblocked vertices and
are connected by unblocked nearest-neighbor paths in the unit
enlargement of that bounding box.  At an edge from $H_C$ to its
complement the exterior vertex has distance five from $C$ and at
least four from every other blocked component.

Choose an exterior-boundary representative $r_H$ for each maximal
hull and define $\rho(z)=r_H$ on $H$, fixing other vertices.
Apply this map \emph{once}.  Each lattice edge has a route between
its two images, using only unblocked vertices and lying in the
original edge together with the enlarged boxes of its incident
hulls.  For completeness, if its endpoints belong to distinct
maximal hulls $H,H'$, the route is
\[
 r_H\longrightarrow v\longrightarrow u\longrightarrow r_{H'}
 \qquad(u\in H,\ v\in H').
\]
The first and last portions run along the respective exterior
boundaries, and the middle portion reverses the original edge.
A representative can lie in another hull, so $\rho$ is not assumed
idempotent.  This is why the once-applied convention matters.
These inputs concern arbitrary finite blocked sets and require no
regular-anchor hypothesis.

Let $n\ge1$ be an integer parameter and assume $10^6dn\le R\le L$.
For a cell $B$, lift its coordinate intervals to $\Z^3$, and take
a coarse rectangular window containing all boxes meeting that lift,
padded by $1000n$ coarse units.  Copy the torus box data periodically
into the window and declare precisely its bad boxes blocked; outside
the window there are no blocked boxes.  Call $B$ \emph{usable} if
every component of this finite blocked set has diameter at most
$2n$.  Lemma~\ref{lem:bad-box-tail} gives
\begin{equation}\label{eq:unusable-cell}
 \PP(B\text{ is unusable})
 \le C_\beta R^3e^{-cn^{\lambda_1}}.
\end{equation}
To verify the periodic issue in this bound, a component of diameter
greater than $2n$ has, from one of its vertices, a path to distance
$n$.  Truncate at its first reach.  The path stays within $n+8$
coarse units of its root; the assumed bound $10^6dn\le L$ ensures
that its projection preserves
these distances and gives a torus bad component reaching $n$.
The window has at most $C_\beta R^3$ vertices, including its periodic
copies, which proves the union bound.  The hull construction therefore
takes place on a finite blocked subset of the covering lattice.

For usable $B$, apply the preceding lattice hull construction.
Each enlarged bounding box has diameter at most $2n+10$; simple
routes in the union of at most two such boxes have length at most
$Cn^3$.  All hulls and routes used near $B$ are strictly inside the
padded window.  There the finite-set designation ``unblocked''
means that the corresponding torus coarse box is good.  Choose
representative cut holes periodically for all good boxes and define
on fine sites of the lifted cell
\begin{equation}\label{eq:comparison-function}
 U_B(x)=u_{i(\rho(z_x))},
\end{equation}
where $z_x$ is the coarse box containing $x$.  Thus $|U_B|\le1$.
This is a function on all sites, even when some sites have no hole.

\begin{lemma}[Comparison across a cell and out of a hull]
\label{lem:hull-comparison}
There is a constant $C_\beta$ such that the following hold for
every usable $B$.
For fine nearest neighbors $x,x'$ within its lift,
\begin{equation}\label{eq:comparison-edge}
 |U_B(x)-U_B(x')|^2
 \le C_\beta n^6\delta^{-1}
       \sum_{\substack{i\in\mathcal H\\
               \dist(x_i,x)\le C_\beta n}} e_i.
\end{equation}
For a cut hole $j$ above $x\in B$, the same right side bounds
$|u_j-U_B(x)|^2$, unless either
\begin{enumerate}[label=(\roman*),nosep]
 \item there is a short-time flag in some box at coarse distance
 at most $100n$ from $z_x$; or
 \item the $\Gamma$-component of $j$ has at most $C_\beta n^3$
 vertices.
\end{enumerate}
The distances in the energy sum may be taken on the torus.
\end{lemma}

\begin{proof}
If $x,x'$ lie in the same coarse box the left side of
\eqref{eq:comparison-edge} vanishes.  Otherwise their box indices
are nearest neighbors.  The imported hull route joins their images
under $\rho$ through good boxes, in a region of diameter $Cn$,
with length at most $Cn^3$.  Replace each of its coarse edges by
the cut-hole paths of Lemma~\ref{lem:good-box-paths}.  This gives
at most $C_\beta n^3$ edges of $\Gamma$ in the stated neighborhood,
each obeying \eqref{eq:contact-cost} with $\tau=\delta$.
Telescoping and Cauchy--Schwarz bound the total squared difference
by path length times the sum of these costs.  Bounding the
multiplicity of any $e_i$ by the path length itself gives the
factor $C_\beta n^6\delta^{-1}$.

Now let $j$ lie over $x$.  If $z_x$ is outside all maximal hulls,
it is good and the local paths already prove the assertion.
Otherwise let $H$ be its maximal hull.  This hull and its enlarged
box have diameter $O(n)$; they project injectively onto the coarse
torus by the same bound $10^6dn\le L$.  The set $J_H$ of cut holes whose
boxes belong to its projection has size at most $C_\beta n^3$.
If the $\Gamma$-component of $j$ stays in $J_H$, exception (ii)
applies.  Otherwise take a simple $\Gamma$-path from $j$ to its
first vertex outside $J_H$.  It has at most $|J_H|$ edges.

To locate these edges, lift a spatial
realization starting at its endpoint in $H$.  If the realization
stays in $H$, it is local.  If it leaves, its first boundary crossing
has an exterior coarse box $z$ with distance at least four from
every bad box.  This crossing lies in the interior of the data
window, so $z$ is genuinely good on the torus.
Lemma~\ref{lem:good-box-paths} confines the \emph{entire}
realization to the $d/2$-neighborhood of that crossing, and its
cut endpoint boxes to coarse distance one from $z$.  In particular
it cannot wind around the torus and return to another copy of $H$.
For every edge on the first-exit path, its realization and both cut
starts thus remain within coarse distance $Cn$ of $z_x$.

It remains to justify its contact costs; the bad boxes inside $H$
did not by themselves guarantee any minimum contact duration.
Choose a contact time on its localized realization, and extend to
the maximal interval on which the two skeleton positions are
fixed.  An endpoint of this interval is $0$, $\beta$, or an actual
mark incident to one of the two occupied slots.  If the interval
had length less than $\delta$, its endpoints would give either a
mark within $\delta$ of an end of the period or two distinct
incident marks within $\delta$ of each other.  The contact sites
are within $12n$ coarse units of $z_x$: the hull diameter is at most
$2n+8$, its boundary boxes are within $2n+9$ of $z_x$, and a
realization meeting such a box stays within one further coarse unit.
Thus $2n+10\le12n$.  Both endpoint marks are incident to these same
fixed positions, including when an endpoint time precedes the chosen
contact time.  Such a flag is
excluded by (i).  Every contact therefore has duration at least
$\delta$.  Endpoint edges need no duration estimate.

For the exiting edge, take its boundary-crossing exterior box $z$
as above.  The final cut hole lies in a box at coarse distance at
most one from $z$.  Clearance makes all boxes within that distance
good.  A coordinate coarse path between them stays in this good
neighborhood, and the local hole paths lead to $i(z)$.  Finally,
$z$ and the representative $r_H$ both lie on the exterior vertex
boundary of $H$.  The imported boundary route joins them through
good boxes in its enlarged bounding box, using at most $Cn^3$
steps after loop erasure.  Lift this route by the local hole paths.
We have joined $j$ to $i(r_H)$, the hole defining $U_B(x)$,
by at most $C_\beta n^3$ controlled edges, all using cut starts
within fine distance $C_\beta n$ of $x$.
The same telescoping estimate proves the assertion.
\end{proof}

Figure~\ref{fig:holes-comparison} illustrates the two local comparisons
used in this argument: energy is charged at cut starts, and an exit
from a hull reaches a nearby surviving boundary route.
\begin{figure}[t]
\centering
\begin{tikzpicture}[x=.92cm,y=.92cm,>=Stealth,
  every node/.style={font=\small},
  hole/.style={circle,fill=black,inner sep=1.7pt}]
  \begin{scope}
    \node[font=\small\bfseries] at (2.4,4.25) {A contact edge};
    \draw[->,gray] (0,3.65)--(5.25,3.65)
      node[pos=.5,above] {one period};
    \draw[gray] (0,3.5)--(0,.3);
    \draw[gray] (5,3.5)--(5,.3);
    \node[above] at (0,3.65) {$0$};
    \node[above] at (5,3.65) {$\beta$};
    \fill[gray!15] (2.1,.4) rectangle (2.85,3.2);
    \draw[thick] (0,2.8)--(1.1,2.8)--(1.1,2.35)
      --(3.3,2.35)--(3.3,2.9)--(5,2.9);
    \draw[thick] (0,.9)--(.8,.9)--(.8,1.35)
      --(3.9,1.35)--(3.9,.9)--(5,.9);
    \node[hole,label=left:$i$] at (0,2.8) {};
    \node[hole,label=left:$j$] at (0,.9) {};
    \node[hole,label=right:$\phi(i)$] at (5,2.9) {};
    \node[hole,label=right:$\phi(j)$] at (5,.9) {};
    \draw[very thick,->] (2.5,2.35)--(2.5,1.35);
    \node[right] at (2.58,1.84) {jump};
    \draw[<->] (2.1,.27)--(2.85,.27)
       node[midway,below] {$\tau\ge\delta$};
    \node[align=center] at (2.6,-.65)
       {$i\longrightarrow\phi(j)\longleftarrow j$\\
        cost uses $e_i+e_j$};
  \end{scope}
  \begin{scope}[shift={(8.0,0)}]
    \node[font=\small\bfseries] at (2.25,4.25)
       {Leaving a filled hull};
    \draw[step=.5,gray!18] (-.25,.2) grid (4.75,3.7);
    \filldraw[fill=gray!12,draw=gray!60,thick]
       (.3,.65)--(3.2,.65)--(3.2,2.7)--(2.45,2.7)
       --(2.45,3.15)--(.3,3.15)--cycle;
    \node at (1,2.7) {$H$};
    \fill[gray!60] (1.2,1.8) rectangle (1.7,2.3);
    \fill[gray!60] (1.7,1.3) rectangle (2.2,1.8);
    \node[font=\scriptsize,align=center] at (1.55,1.03)
      {bad boxes};
    \node[hole,label=left:$j$] (j) at (.7,1.8) {};
    \node[hole] (v1) at (1.2,1.45) {};
    \node[hole] (v2) at (2.5,1.5) {};
    \node[hole] (v3) at (3.45,1.85) {};
    \draw[thick,->] (j)--(v1);
    \draw[thick,->] (v1)--(v2);
    \draw[very thick,->] (v2)--(v3);
    \draw[rounded corners,thick,->]
      (v3)--(3.65,1.85)--(3.65,3.45)--(2.9,3.45);
    \node[hole,label=above:$i(r_H)$] at (2.9,3.45) {};
    \draw[gray,dashed] (3.08,1.4) rectangle (3.92,2.3);
    \node[font=\scriptsize,align=left,anchor=west]
      at (3.96,1.85) {good\\boundary};
    \node[align=center] at (2.25,-.65)
      {first exit, then an exterior route\\
       to the comparison value};
  \end{scope}
\end{tikzpicture}
\caption{Two parts of the hole comparison, drawn schematically.
Left: a fresh walk jumps between skeleton lines during a contact
interval.  Comparing both cut values through $u_{\phi(j)}$ charges
energy only to the cut starts $i,j$.
Right: a path in $\Gamma$ first exits a filled coarse hull.
Clearance makes its crossing box good; the absence of long mark paths
then confines the whole exiting link near that box.  A surviving
exterior-boundary route leads to $i(r_H)$, which defines $U_B$ on
the hull.  Gray squares indicate bad boxes; the larger gray region
represents the filled hull, with distances and thicknesses not to scale.}
\label{fig:holes-comparison}
\end{figure}

\FloatBarrier

\subsection{Energy, exceptional sets, and the final estimate}

The construction has supplied a bounded comparison function on each
usable cell.  Its lattice gradients are controlled by the total
traversal energy, and a hole fails to match that function only near a
short-time flag or in a small closed class.  We now sum these three
contributions and choose the auxiliary scales.

\begin{lemma}[Cellwise approximation of the resolvent]
\label{lem:cell-resolvent}
There are random numbers $c_B\in[-1,1]$, one per cell, such that
\begin{equation}\label{eq:cell-resolvent}
 \E\sum_B\sum_{\substack{i\in\mathcal H\\x_i\in B}}
        (u_i-c_B)^2\le C_\beta N R^{-1/20}.
\end{equation}
\end{lemma}

\begin{proof}
On a usable cell let $c_B=|B|^{-1}\sum_{x\in B}U_B(x)$; on an
unusable cell put $c_B=0$.  The rectangular lattice Poincar\'e
inequality, whose constant is at most $CR^2$ for the stated side
lengths, gives
\[
 \sum_{x\in B}|U_B(x)-c_B|^2
 \le CR^2\sum_{\substack{\{x,x'\}\subset B\\|x-x'|_1=1}}
                  |U_B(x)-U_B(x')|^2.
\]
It holds in the chosen lift, also for a cell covering an entire
coordinate circle: only the internal rectangle edges are needed.
Use Lemma~\ref{lem:hull-comparison} both here and for the
nonexceptional differences $u_i-U_B(x_i)$.  There are at most
$\ell$ holes above a site.  In summing all the local energy
sums, any fixed torus $e_i$ occurs at most $C_\beta n^3$ times,
since the original fine sites, over all cells, partition the torus.
Periodic copies used to construct the windows do not add summation
sites.  Equation~\eqref{eq:traversal-energy} therefore bounds these
contributions by
\[
 C_\beta R^2n^9\delta^{-1}qN.
\]

All remaining squared differences are at most four.  A fixed fine
site has at most $Cn^3$ boxes in its short-flag neighborhood; its
exceptional probability is at most $C_\beta n^3\delta$ by
Lemma~\ref{lem:bad-box-tail}.  The total number of holes in
$\Gamma$-components of size at most $C_\beta n^3$ is at most
$C_\beta n^3$ times the number of components, whose expectation
is bounded by \eqref{eq:component-count}.  Finally sum
\eqref{eq:unusable-cell} with weights at most $\ell|B|$.
Together these estimates give
\begin{equation}\label{eq:comparison-errors}
 \frac1N\E\sum_B\sum_{\substack{i\in\mathcal H\\x_i\in B}}
      (u_i-c_B)^2
 \le C_\beta\left(
 R^2n^9\delta^{-1}q+n^3\delta+n^3q
                   +R^3e^{-cn^{\lambda_1}}\right).
\end{equation}

Take now
\[
 \delta=R^{-1/8},\qquad n=\lceil R^{1/200}\rceil.
\]
At fixed $\beta$, sufficiently large $R$ meets all preceding
smallness conditions, including $10^6dn\le R\le L$, and the
short-flag threshold in Lemma~\ref{lem:bad-box-tail}.  The first
three terms on the right of \eqref{eq:comparison-errors} are,
respectively,
\[
 O(R^{-33/100}),\qquad O(R^{-11/100}),\qquad
 O(R^{-497/200}).
\]
The fourth is smaller than any negative power of $R$.  This proves
\eqref{eq:cell-resolvent} with the stated, nonoptimal exponent.
\end{proof}

\begin{proof}[Proof of Proposition~\ref{prop:weak-field}]
The partition has at most $CV/R^3$ cells and each has volume
comparable to $R^3$.  Lemmas~\ref{lem:cell-resolvent}
and~\ref{lem:hole-cancellation}, together with $|c_B|\le1$,
give
\begin{align*}
 \E\bigl|\ip{f_{\mathcal H}}{u}\bigr|
 &\le \E\sum_B\sum_{i\in\mathcal H,\,x_i\in B}|u_i-c_B|
       +\sum_B\E\left|\sum_{i\in\mathcal H,\,x_i\in B}f_i\right|\\
 &\le C_\beta N R^{-1/40}
          +CV(qR^3)^{-1/2}.
\end{align*}
In \eqref{eq:transverse-resolvent}, add the diagonal terms inside
the resolvent sum and discard their nonnegative subtraction.
Since $u=qD^{-1}f_{\mathcal H}$, this yields
\[
 \frac qV\langle D_f^2\rangle_b
 \le \frac{qN}{4V}+
       \frac1{2V}\E\bigl|\ip{f_{\mathcal H}}{u}\bigr|
 \le C_\beta\bigl(q+R^{-1/40}+R^{-1/4}\bigr).
\]
This is \eqref{eq:weak-field}.  The assumptions needed for the
sparse-pin estimate hold because $q=R^{-5/2}\ge L^{-5/2}$ and
the lower bound on $R$ can include $L_*(\beta)$.
\end{proof}
